% Capitulo: Antes e depois da IA no meu trabalho
\chapter{Antes e depois da IA no meu trabalho}
\label{cap:antes-depois}

\rascunho{
  Capítulo escrito como síntese das mudanças discutidas no livro. Acrescentar, em cada seção,
  um exemplo concreto do meu trabalho (sem nomes e com números só em ordem de grandeza). Se
  preferir um título menos pessoal: \enquote{Antes e depois da IA no trabalho com dados}.
}

Ao longo do livro, a IA apareceu em quase todos os capítulos: mudando o que vale estudar, como
os dados precisam ser organizados, como o código é escrito e revisado, quanto as coisas custam
e como os times se estruturam. Este capítulo junta essas mudanças num lugar só e compara, tema a
tema, o mesmo fluxo de trabalho de um time de dados antes e depois da IA. Não é uma lista de
vantagens. Algumas coisas melhoraram muito, outras pioraram, e quase todas exigiram cuidados
novos.

\section{Estimativa e refinamento}
\label{sec:ad-estimativa}

A parte do trabalho que mais encolheu foi a construção: escrever a consulta, montar o modelo,
gerar a primeira versão do painel. Isso mudou o tamanho relativo das tarefas. O que antes
ocupava a maior parte de uma história agora ocupa uma fração, e o que antes era rápido, como
entender a pergunta de negócio e investigar os dados, passou a dominar o tempo.

\antesdepois{A estimativa era dominada pela construção. Uma história grande era uma história
com muito código para escrever, e o refinamento discutia principalmente como fazer.}
{A construção encolheu e a descoberta passou a dominar. O refinamento discute principalmente o
que fazer, quais são as regras e como saber se está certo.}

A consequência para o refinamento é que ele ficou mais importante, e não menos. Uma história mal
definida, que antes era descoberta no meio da construção, agora é construída rapidamente e da
forma errada. A separação entre descoberta e entrega (\veja{sec:descoberta}) e os critérios de
entrada no quadro (\veja{sec:refinamento}) passaram a ser a principal defesa contra retrabalho.

\rascunho{Exemplo meu: como mudou o tamanho ou a natureza das histórias no meu time.}

\section{Código e revisão}
\label{sec:ad-codigo}

Escrever código ficou muito mais rápido. Revisar ficou mais importante, e não ficou mais fácil.
O volume de código produzido cresceu, e boa parte dele foi escrita por alguém, a IA, que não
estará na reunião quando o número estiver errado (capítulo~\ref{cap:programar}).

\antesdepois{O autor escrevia devagar e entendia cada linha. A revisão conferia principalmente
estilo e regras de negócio.}
{O código chega rápido e em volume. A revisão passou a ser o principal controle de qualidade, e
precisa conferir também se o autor entendeu o que a IA escreveu.}

Na prática, isso levou a algumas regras que antes pareciam excesso de zelo: alterações pequenas,
mesmo quando a IA consegue gerar uma grande; o autor revisando o próprio código gerado antes de
pedir revisão; e as perguntas de grão e de junção feitas em toda revisão de SQL
(\veja{sec:revisao-codigo}).

\rascunho{Exemplo meu: um caso em que a revisão pegou um erro de código gerado por IA.}

\section{Testes}
\label{sec:ad-testes}

Talvez a mudança mais importante tenha sido nos testes. O desenvolvimento guiado por testes
sempre foi recomendado e quase nunca praticado em dados, porque escrever testes custava um tempo
que o prazo não dava. Com a IA, esse custo caiu, e os testes passaram a vir antes dos modelos
(\veja{sec:tdd}).

\antesdepois{Os testes ficavam para quando sobrasse tempo. As regras de negócio eram conferidas
com o stakeholder a cada versão do modelo, o que exigia a presença dele em cada iteração.}
{Os testes vêm antes do modelo e viram a especificação da regra. A validação com o stakeholder
acontece uma vez, sobre os casos de teste, e o time itera sem depender da agenda dele.}

\rascunho{Exemplo meu: um projeto em que escrever os testes antes mudou o tempo de validação.}

\section{Documentação}
\label{sec:ad-documentacao}

A documentação era a eterna pendência de todo time de dados: todos concordavam que era
importante, e ela sempre ficava para depois. Duas coisas mudaram. A IA passou a escrever a
primeira versão das descrições de tabelas, colunas e métricas, a partir do código e dos dados. E a
documentação ganhou um consumidor novo, os assistentes que respondem perguntas sobre os dados, para
os quais ela é insumo essencial (capítulo~\ref{cap:contexto-ia}).

\antesdepois{A documentação ficava num documento à parte, desatualizado, e servia
principalmente para quem chegava ao time.}
{A documentação vive junto com o código, a IA escreve a primeira versão e alguém que conhece o
negócio revisa. Ela serve ao time e à IA, que a usa a cada pergunta.}

O mesmo vale para a transferência de conhecimento. Explicar um modelo antigo a uma pessoa nova
ficou muito mais fácil quando um assistente consegue ler o código e responder às perguntas dela.
O que não ficou mais fácil é transferir o conhecimento de negócio, que continua dependendo de
pessoas.

\section{Custo}
\label{sec:ad-custo}

A forma de pagar pelas ferramentas mudou junto com elas. O que era uma licença fixa por usuário
passou a ser, cada vez mais, consumo: créditos de processamento no armazém de dados e tokens nos
modelos de linguagem (capítulo~\ref{cap:custo-consumo}).

\antesdepois{O custo das ferramentas era fixo e previsível, negociado na renovação. O gerente
controlava o acesso.}
{O custo acompanha o uso. O gerente acompanha o consumo de forma contínua, por time e por
produto, e a eficiência técnica aparece diretamente na fatura.}

\rascunho{Exemplo meu: como passei a acompanhar consumo, e uma otimização que fez diferença.}

\section{O que exigiu novos cuidados}
\label{sec:ad-cuidados}

Nem tudo melhorou. Algumas mudanças exigiram cuidados que antes não existiam, e vale
reuni-los aqui.

O primeiro é a privacidade. Ficou fácil colar num assistente um trecho de dados que nunca
deveria sair da empresa, e as regras sobre o que pode e o que não pode ir para um modelo
precisaram ser escritas e comunicadas (\veja{sec:seguranca}).

O segundo é a confiança excessiva. O código gerado parece correto, a resposta vem com
segurança, e a tendência humana é aceitar. Os estudos citados ao longo do livro mostram que
essa confiança nem sempre é justificada, e que a percepção de ganho de produtividade pode não
corresponder ao ganho real \parencite{becker2025,perry2023}.

O terceiro é a formação de quem está chegando. Quem aprende com a IA fazendo tudo aprende menos
\parencite{shen2026}, e um time que só contrata quem já sabe vai ficar sem seniores no futuro
(\veja{sec:ia-comprime}).

O quarto é o custo invisível do excesso. Ficou tão fácil criar painéis, modelos, aplicações e
agentes que a quantidade cresce mais rápido do que a capacidade de manter. As métricas de uso
dos próprios produtos de dados (\veja{sec:metricas-uso}) e a disciplina de desligar o que não é
usado passaram a ser parte do trabalho.

\begin{destaque}
A IA não diminuiu o trabalho de um time de dados. Ela mudou de lugar: menos construção, mais
definição, revisão, governança e conhecimento de negócio.
\end{destaque}

\rascunho{Exemplo meu: o cuidado novo que mais me custou aprender.}

\begin{leituras}
  \leitura{dellacqua2023}{Para entender por que a IA ajuda muito em algumas tarefas e
    atrapalha em outras, e como identificar a diferença.}
  \leitura{becker2025}{Para a distância entre o ganho percebido e o ganho medido, um bom
    antídoto contra o entusiasmo sem critério.}
  \leitura{forsgren2018}{Para medir, com indicadores objetivos, se a forma de trabalhar do
    time de fato melhorou.}
\end{leituras}
